\documentclass{article}
% Encodage
\usepackage[utf8]{inputenc}
\usepackage[T1]{fontenc}
\usepackage[french]{babel}
% Mise en page
\usepackage{geometry}
\geometry{margin=2cm}
% Package Maths
\usepackage{amsmath}
% Package unités SI
\usepackage{siunitx}
\DeclareSIUnit{\gauss}{G}
\DeclareSIUnit{\tesla}{T}
% Package Hyperlien
\usepackage{hyperref}
\title{Projet Intégratif : Itinéraire}
\author{Adam EL MARZOUKI \and Noémie MORNON \and Théo SAVARY \and Roger-Parfait MBAITORNDIA \and Axel LÊ TRONG}
\date{\today}
\begin{document}
\maketitle
\tableofcontents
\newpage
\section{Introduction}
\section{Partie 1}
\subsection{Objectif}
L'objectif de cette partie est de développer un système autonome d'acquisition et de traitement de données permettant d'estimer la trajectoire d'un utilisateur à partir de plusieurs capteurs embarqués sur l'ITIPhone. Nous avions plusieurs capteurs à disposition. Afin de concevoir un algorithme de navigation adapté au cahier des charges, nous avons choisi d'utiliser uniquement le GPS et l'IMU.
De plus, nous avons fait l'hypothèse que l'ITIPhone reste positionné à plat tout au long du trajet. Les rotations autour de l'axe y et x sont donc négligeables. La dynamique angulaire est réduite à une rotation autour de l'axe vertical $\mathbf{z}$. 
Bien que le GPS fournisse une position correcte extérieur, il reste sensible aux perturbations environnementales et peut produire des mesures bruitées ou aberrantes. L'idée est donc de combiner les informations du GPS avec celles fournies par l'IMU à l'aide du filtre de Kalman . 

\subsection{Analyse préliminaire des données GPS}
Avant la mise en place du filtre, nous avons étudié les données GNSS obtenues lors des différentes acquisitions. Nous avons observé qu'avec une météo ensoleillée, le GPS fournissait déja des positions relativement précises. En revanche, certaines mesures avaient des sauts importants dus à des erreurs de réception. Ces erreurs étaient importantes lorsqu'on se trouvait à proximité d'obstacles ou à l'intérieur d'un batiment. Afin de limiter leur impact nous avons mis en place une étape de nettoyage des mesures GNSS

\subsection{Suppression des points aberrants}
Pour détecter les mesures aberrantes, nous avons calculé pour chaque point : 
\begin{itemize}
    \item la distance au point précédent;
    \item la distance au point suivant;
    \item la distance entre les deux points
\end{itemize}

Un point est considéré comme aberrant lorsqu'il est éloigné de ses deux voisins, mais les deux voisins eux sont cohérents entre eux. Cette étape nous permet donc d'éliminer les sauts ponctuels du GPS avant leur utilisation.

\subsection{Changement de repère et projection}
Les coordonnées GNSS sont exprimées en coordonnées sphériques sous la forme latitude/longitude. Ce format est adapté à l'affichage cartographique mais peut adapté pour les calculs. Nous avons donc converti les positions dans un repère cartésien exprimé en mètres. Le premier point de l'acquisition est utilisé comme origine du repère.  Les coordonnées locales sont obtenues par : 
\[
\begin{aligned}
x &= \text{rad}(\text{lon}-\text{lon}_0) \cdot R_{\text{earth}} \cdot \cos(\text{rad}(\text{lat}_0)) \\
y &= \text{rad}(\text{lat}-\text{lat}_0) \cdot R_{\text{earth}}
\end{aligned}
\]
où $R_{\text{earth}} = \SI{6371000}{m}$ représente le rayon moyen de la Terre et $\text{rad}(\cdot)$ la conversion en radians.

\subsection{Calibration et estimation des biais}

Lors des premiers essais, nous avons rapidement constaté que l'utilisation directe des mesures de l'IMU provoquait une dérive importante de la trajectoire estimée. Pour corriger ce défaut, nous avons réaliser plusieurs acquisitions en laissant le système immobile. Dans ces conditions, les accélérations et vitesses angulaires mesurées devraient théoriquement être nulles ( ou égales à la gravité ). Les écarts observées correspondent donc aux biais des capteurs. Nous avons estimé ces biais à partir de la moyenne empirique des acquisitions immobiles enregistrées : 

\[
b_{a,x}=\frac1N\sum_{k=1}^{N}a_{x,k}
\]

\[
b_{a,y}=\frac1N\sum_{k=1}^{N}a_{y,k}
\]

\[
b_{g,z}=\frac1N\sum_{k=1}^{N}\omega_{z,k}
\]

Ces valeurs sont ensuite retirées aux mesures avant leur utilisation dans le filtre. De plus, la variance de ces signaux statiques nous permet de dimensionner les matrices de covariance du bruit du filtre.

\subsection{Estimation du cap initial}
Conformément à notre hypothèse de travail (téléphone positionné à plat), seul le cap horizontal $\theta$ doit être initialisé. Nous exploitons le magnétomètre embarqué qui mesure les composantes $m_x$ et $m_y$ du champ magnétique terrestre. À l'arrêt, le cap initial $\theta_0$ est déterminé par :
\[
\theta_0 = \arctan2(m_y, m_x)
\]

Cette valeur permet d'orienter correctement le vecteur d'état initial avant le début du mouvement. 

\subsection{Filtre de Kalman Etendu (EKF)}
L'objectif est de combiner les données du GPS avec celle de l'IMU tout en profitant des avantages de chacune. Les données de l'IMU sont utilisées pour prédire l'évolution de la trajectoire entre deux mesures GPS succesives. La projection des accélérations du repère du téléphone vers le repère terrestre dépend de l'orientation estimée du système. Cette opération introduit des termes trigonométriques dans le modèle, ce qui rend la dynamique du système non linéaire. Dans ce contexte le filtre de Kalman n'est pas idéal. Nous utilisons un filtre de Kalman Etendu (EKF) capable de prendre en compte ces non-linéarités. 

\subsubsection{Vecteur d'état}

L'état du système est défini par :

\[
\mathbf{x}_k = \begin{pmatrix} x_k \\ y_k \\ v_{x,k} \\ v_{y,k} \\ \theta_k \end{pmatrix}
\]
où $x_k, y_k$ sont les positions en mètres, $v_{x,k}, v_{y,k}$ les vitesses linéaires dans le repère monde, et $\theta_k$ le cap.

\subsubsection{Étape de prédiction}
À chaque pas de temps $\Delta t$, les accélérations mesurées dans le repère du téléphone sont corrigées de leurs biais respectifs puis projetées dans le repère global (monde) via la matrice de rotation plane 2D :
\[
\begin{pmatrix} a_{x,k}^{\text{monde}} \\ a_{y,k}^{\text{monde}} \end{pmatrix} = 
\begin{pmatrix} \cos\theta_{k-1} & -\sin\theta_{k-1} \\ \sin\theta_{k-1} & \cos\theta_{k-1} \end{pmatrix}
\begin{pmatrix} a_{x,k} - b_{a,x} \\ a_{y,k} - b_{a,y} \end{pmatrix}
\]

L'état prédit $\mathbf{x}_{k|k-1}$ est mis à jour à l'aide des équations de la cinématique :
\[
\mathbf{x}_{k|k-1} = \begin{pmatrix}
x_{k-1} + v_{x,k-1}\Delta t + \frac{1}{2}a_{x,k}^{\text{monde}}\Delta t^2 \\
y_{k-1} + v_{y,k-1}\Delta t + \frac{1}{2}a_{y,k}^{\text{monde}}\Delta t^2 \\
v_{x,k-1} + a_{x,k}^{\text{monde}}\Delta t \\
v_{y,k-1} + a_{y,k}^{\text{monde}}\Delta t \\
\theta_{k-1} + (\omega_{z,k} - b_{g,z})\Delta t
\end{pmatrix}
\]

\subsubsection{Propagation de l'incertitude}

Le filtre ne propage pas uniquement un état mais également l'incertitude associée à cet état.

Pour cela, il utilise la matrice de transition :

\[
\mathbf{F} = \begin{pmatrix}
1 & 0 & \Delta t & 0 & 0 \\
0 & 1 & 0 & \Delta t & 0 \\
0 & 0 & 1 & 0 & 0 \\
0 & 0 & 0 & 1 & 0 \\
0 & 0 & 0 & 0 & 1
\end{pmatrix}
\]

Cette matrice décrit l'influence de l'état courant sur l'état futur.

La covariance est alors propagée par :
\[
\mathbf{P}_{k|k-1} = \mathbf{F} \mathbf{P}_{k-1|k-1} \mathbf{F}^\top + \mathbf{Q}
\]
où $\mathbf{Q}$ désigne la matrice de covariance du bruit de processus.

\subsubsection{Correction GNSS}

Le récepteur GNSS mesurant uniquement la position bidimensionnelle, la matrice d'observation s'écrit :
\[
\mathbf{H} = \begin{pmatrix} 1 & 0 & 0 & 0 & 0 \\ 0 & 1 & 0 & 0 & 0 \end{pmatrix}
\]

Afin de tenir compte de la qualité variable du signal GNSS, nous utilisons l'indicateur HDOP fourni par le récepteur.

Lorsque le HDOP augmente, la confiance accordée aux mesures GPS diminue. Cette information est intégrée directement dans la matrice de covariance des mesures :

\[
R_k = R_{base}\times1.5\times HDOP_k^3
\]

\subsubsection{Innovation et covariance de l'innovation}
L'innovation $\mathbf r_k$ est la différence entre notre mesure $\mathbf y_k$ et notre mesure prédite $\mathbf h(x_{k|k-1})$, autrement dit elle représente la «~suprise~». On la définit ainsi :
\[
\mathbf r_k = \mathbf y_k - \mathbf h(x_{k|k-1})
\]
De plus, la covariance de l'innovation est calculée itérativement de la manière suivante :
\[
\begin{aligned}
    \mathbf S_k &= \mathbf H_k \, \mathbf P_{k|k-1} \, \mathbf H_k^\top + \mathbf R \\
                &= \mathbf H \, \mathbf P_{k|k-1} \, \mathbf H^\top + \mathbf R
    \quad \text{car } \mathbf H_k = \mathbf H
\end{aligned}
\]
avec comme cas initial :
\[
\mathbf S_0 = \mathbf H \, \mathbf P_{0|0} \, \mathbf H^\top + \mathbf R_0
\]

\subsubsection{Gain de Kalman K}
Le gain de Kalman $\mathbf K_k$ est une matrice qui représente, à l'itération $k$, si l'on devrait plutôt faire confiance aux mesures de notre capteur GNSS ou aux prédictions de notre modèle. On l'exprime ainsi :
\[
\begin{aligned}
    \mathbf K_k &= \mathbf P_{k|k-1} \, \mathbf H_k^\top \, \mathbf S_k^{-1} \\
                &= \mathbf P_{k|k-1} \, \mathbf H^\top \, \mathbf S_k^{-1}
    \quad \text{car } \mathbf H_k = \mathbf H
\end{aligned}
\]
avec comme cas initial :
\[
\mathbf K_0 = \mathbf P_{0|0} \, \mathbf H^\top \, \mathbf S_0^{-1}
\]

\subsubsection{Correction}
Une fois le gain de Kalman calculé il ne nous reste plus qu'à corriger notre prédiction de cette manière :
\[
\begin{aligned}
    \mathbf x_{k|k} &= \mathbf x_{k|k-1} + \mathbf K_k \, \mathbf r_k, \\
    \mathbf P_{k|k} &= (\mathbf I - \mathbf K_k \, \mathbf H_k) \, \mathbf P_{k|k-1}
\end{aligned}
\]


\section{Partie 2}
\section{Difficultées rencontrées}
\subsection{Partie 1}
\subsection{Partie 2}
\section{Répartition du travail}
\end{document}
